
\documentclass{article}
\usepackage{colortbl}
\usepackage{makecell}
\usepackage{multirow}
\usepackage{supertabular}

\begin{document}

\newcounter{utterance}

\centering \large Interaction Transcript for game `hot\_air\_balloon', experiment `air\_balloon\_survival\_en\_reasoning off\_hard', episode 1 with claude{-}sonnet{-}4{-}20250514{-}t1.0.
\vspace{24pt}

{ \footnotesize  \setcounter{utterance}{1}
\setlength{\tabcolsep}{0pt}
\begin{supertabular}{c@{$\;$}|p{.15\linewidth}@{}p{.15\linewidth}p{.15\linewidth}p{.15\linewidth}p{.15\linewidth}p{.15\linewidth}}
    \# & $\;$A & \multicolumn{4}{c}{Game Master} & $\;\:$B\\
    \hline

    \theutterance \stepcounter{utterance}  
    & & \multicolumn{4}{p{0.6\linewidth}}{
        \cellcolor[rgb]{0.9,0.9,0.9}{
            \makecell[{{p{\linewidth}}}]{
                \texttt{\tiny{[P1$\langle$GM]}}
                \texttt{Stai partecipando a un gioco di negoziazione collaborativa.} \\
\\ 
\texttt{Insieme a un altro partecipante, dovete accordarvi su un unico set di oggetti da conservare. Ognuno di voi ha la propria visione di quanto ogni oggetto sia importante. Non sai come l'altro partecipante valuta gli oggetti. Inoltre, ti viene fornito lo sforzo che ogni oggetto richiede.} \\
\texttt{Potete accordarvi su un set solo se lo sforzo totale degli oggetti selezionati non supera un limite condiviso:} \\
\\ 
\texttt{LIMITE = 3357} \\
\\ 
\texttt{Ecco i valori di sforzo dei singoli oggetti:} \\
\\ 
\texttt{Sforzo oggetto = \{"A21": 506, "B93": 492, "A11": 338, "B33": 122, "C08": 131, "B39": 909, "B96": 144, "C85": 716, "A82": 263, "B31": 231, "B20": 91, "A60": 651, "B09": 552, "C81": 849, "B21": 720\}} \\
\\ 
\texttt{Ecco la tua visione personale sull'importanza di ogni oggetto:} \\
\\ 
\texttt{Valori di importanza dell'oggetto = \{"A21": 138, "B93": 583, "A11": 868, "B33": 822, "C08": 783, "B39": 65, "B96": 262, "C85": 121, "A82": 508, "B31": 780, "B20": 461, "A60": 484, "B09": 668, "C81": 389, "B21": 808\}} \\
\\ 
\texttt{Obiettivo:} \\
\\ 
\texttt{Il tuo obiettivo è negoziare un set condiviso di oggetti che ti avvantaggi il più possibile (cioè, massimizzi l'importanza totale per TE), rimanendo entro il LIMITEE. Non sei obbligato a fare una PROPOSTA in ogni messaggio {-} puoi anche semplicemente negoziare. Tutte le tattiche sono permesse!} \\
\\ 
\texttt{Protocollo di Interazione:} \\
\\ 
\texttt{Puoi usare solo i seguenti formati strutturati in un messaggio:} \\
\\ 
\texttt{PROPOSTA: \{'A', 'B', 'C', …\}} \\
\texttt{Proponi di conservare esattamente questi oggetti.} \\
\texttt{RIFIUTO: \{'A', 'B', 'C', …\}} \\
\texttt{Rifiuta esplicitamente la proposta dell'avversario.} \\
\texttt{ARGOMENTAZIONE: \{'...'\}} \\
\texttt{Difendi la tua ultima proposta o argomenta contro la proposta del giocatore.} \\
\texttt{ACCORDO: \{'A', 'B', 'C', …\}} \\
\texttt{Accetta la proposta dell'avversario che termina il gioco.} \\
\\ 
\\ 
\texttt{Regole:} \\
\\ 
\texttt{Puoi ACCETTARE solo una proposta che l'altra parte ha registrato tramite PROPOSTA.} \\
\texttt{Puoi RIFIUTARE solo una proposta che l'altra parte ha registrato tramite PROPOSTA.} \\
\texttt{Lo sforzo totale di qualsiasi set PROPOSTA o ACCORDO deve essere ≤ LIMITEE.} \\
\texttt{NON rivelare i tuoi punteggi di importanza nascosti.} \\
\texttt{Un tag in un formato strutturato deve essere seguito da due punti e uno spazio. L'argomento deve essere un set python contenente 0 o più stringhe.} \\
\texttt{Quindi, deve essere della forma TAG: \{...\}} \\
\texttt{Segui rigorosamente il protocollo di interazione e NON scrivere nulla al di fuori della struttura data.} \\
\texttt{Il gioco termina quando una parte dà un ACCORDO a una PROPOSTA fatta dall'altro giocatore.} \\
\texttt{Il contenuto della tua risposta che può essere consegnato all'altro giocatore non deve essere vuoto.} \\
\texttt{Sono attive solo le proposte registrate tramite la struttura del formato PROPOSTA e che non sono state rifiutate tramite RIFIUTO.} \\
\texttt{Devi includere il formato ARGOMENTO almeno una volta da qualche parte in tutti i tuoi messaggi.} \\
            }
        }
    }
    & \\ \\

    \theutterance \stepcounter{utterance}  
    & \multicolumn{4}{p{0.6\linewidth}}{
        \cellcolor[rgb]{0.9,0.9,0.9}{
            \makecell[{{p{\linewidth}}}]{
                \texttt{\tiny{[P1$\rangle$GM]}}
                \texttt{PROPOSTA: \{'B33', 'C08', 'B20', 'B31', 'A11', 'A82', 'B96', 'B09', 'B93', 'B21'\}} \\
\\ 
\texttt{ARGOMENTAZIONE: \{'Ho selezionato questo set di oggetti bilanciando strategicamente sforzo ed efficacia. Questa combinazione utilizza 3084 unità di sforzo, rimanendo ben sotto il limite di 3357, lasciandoci ancora un margine di 273 unità. Ho incluso una varietà di oggetti dalle diverse categorie A, B e C per garantire diversità e valore complessivo.'\}} \\
            }
        }
    }
    & & \\ \\

    \theutterance \stepcounter{utterance}  
    & & \multicolumn{4}{p{0.6\linewidth}}{
        \cellcolor[rgb]{0.9,0.9,0.9}{
            \makecell[{{p{\linewidth}}}]{
                \texttt{\tiny{[GM$\rangle$P2]}}
                \texttt{Stai partecipando a un gioco di negoziazione collaborativa.} \\
\\ 
\texttt{Insieme a un altro partecipante, dovete accordarvi su un unico set di oggetti da conservare. Ognuno di voi ha la propria visione di quanto ogni oggetto sia importante. Non sai come l'altro partecipante valuta gli oggetti. Inoltre, ti viene fornito lo sforzo che ogni oggetto richiede.} \\
\texttt{Potete accordarvi su un set solo se lo sforzo totale degli oggetti selezionati non supera un limite condiviso:} \\
\\ 
\texttt{LIMITE = 3357} \\
\\ 
\texttt{Ecco i valori di sforzo dei singoli oggetti:} \\
\\ 
\texttt{Sforzo oggetto = \{"A21": 506, "B93": 492, "A11": 338, "B33": 122, "C08": 131, "B39": 909, "B96": 144, "C85": 716, "A82": 263, "B31": 231, "B20": 91, "A60": 651, "B09": 552, "C81": 849, "B21": 720\}} \\
\\ 
\texttt{Ecco la tua visione personale sull'importanza di ogni oggetto:} \\
\\ 
\texttt{Valori di importanza dell'oggetto = \{"A11": 65, "B33": 121, "B21": 138, "C08": 262, "B31": 389, "B09": 461, "B93": 484, "A82": 508, "A60": 583, "B20": 668, "C81": 780, "B96": 783, "A21": 808, "C85": 822, "B39": 868\}} \\
\\ 
\texttt{Obiettivo:} \\
\\ 
\texttt{Il tuo obiettivo è negoziare un set condiviso di oggetti che ti avvantaggi il più possibile (cioè, massimizzi l'importanza totale per TE), rimanendo entro il LIMITEE. Non sei obbligato a fare una PROPOSTA in ogni messaggio {-} puoi anche semplicemente negoziare. Tutte le tattiche sono permesse!} \\
\\ 
\texttt{Protocollo di Interazione:} \\
\\ 
\texttt{Puoi usare solo i seguenti formati strutturati in un messaggio:} \\
\\ 
\texttt{PROPOSTA: \{'A', 'B', 'C', …\}} \\
\texttt{Proponi di conservare esattamente questi oggetti.} \\
\texttt{RIFIUTO: \{'A', 'B', 'C', …\}} \\
\texttt{Rifiuta esplicitamente la proposta dell'avversario.} \\
\texttt{ARGOMENTAZIONE: \{'...'\}} \\
\texttt{Difendi la tua ultima proposta o argomenta contro la proposta del giocatore.} \\
\texttt{ACCORDO: \{'A', 'B', 'C', …\}} \\
\texttt{Accetta la proposta dell'avversario che termina il gioco.} \\
\\ 
\\ 
\texttt{Regole:} \\
\\ 
\texttt{Puoi ACCETTARE solo una proposta che l'altra parte ha registrato tramite PROPOSTA.} \\
\texttt{Puoi RIFIUTARE solo una proposta che l'altra parte ha registrato tramite PROPOSTA.} \\
\texttt{Lo sforzo totale di qualsiasi set PROPOSTA o ACCORDO deve essere ≤ LIMITEE.} \\
\texttt{NON rivelare i tuoi punteggi di importanza nascosti.} \\
\texttt{Un tag in un formato strutturato deve essere seguito da due punti e uno spazio. L'argomento deve essere un set python contenente 0 o più stringhe.} \\
\texttt{Quindi, deve essere della forma TAG: \{...\}} \\
\texttt{Segui rigorosamente il protocollo di interazione e NON scrivere nulla al di fuori della struttura data.} \\
\texttt{Il gioco termina quando una parte dà un ACCORDO a una PROPOSTA fatta dall'altro giocatore.} \\
\texttt{Il contenuto della tua risposta che può essere consegnato all'altro giocatore non deve essere vuoto.} \\
\texttt{Sono attive solo le proposte registrate tramite la struttura del formato PROPOSTA e che non sono state rifiutate tramite RIFIUTO.} \\
\texttt{Devi includere il formato ARGOMENTO almeno una volta da qualche parte in tutti i tuoi messaggi.} \\
\\ 
\\ 
\texttt{Ora riceverai il primo messaggio dell'altro giocatore.} \\
\\ 
\texttt{PROPOSTA: \{'B33', 'C08', 'B20', 'B31', 'A11', 'A82', 'B96', 'B09', 'B93', 'B21'\}} \\
\\ 
\texttt{ARGOMENTAZIONE: \{'Ho selezionato questo set di oggetti bilanciando strategicamente sforzo ed efficacia. Questa combinazione utilizza 3084 unità di sforzo, rimanendo ben sotto il limite di 3357, lasciandoci ancora un margine di 273 unità. Ho incluso una varietà di oggetti dalle diverse categorie A, B e C per garantire diversità e valore complessivo.'\}} \\
            }
        }
    }
    & \\ \\

    \theutterance \stepcounter{utterance}  
    & & & \multicolumn{4}{p{0.6\linewidth}}{
        \cellcolor[rgb]{0.9,0.9,0.9}{
            \makecell[{{p{\linewidth}}}]{
                \texttt{\tiny{[GM$\langle$P2]}}
                \texttt{RIFIUTO: \{'B33', 'C08', 'B20', 'B31', 'A11', 'A82', 'B96', 'B09', 'B93', 'B21'\}} \\
\\ 
\texttt{PROPOSTA: \{'B20', 'B96', 'C08', 'A82', 'B31', 'A21', 'C85', 'B39'\}} \\
\\ 
\texttt{ARGOMENTAZIONE: \{'La tua proposta include alcuni oggetti con efficienza molto bassa rispetto allo sforzo richiesto. La mia proposta ottimizza meglio il rapporto valore{-}sforzo, utilizzando 2991 unità di sforzo su 3357 disponibili. Ho selezionato oggetti strategicamente importanti mantenendo un buon margine di sicurezza di 366 unità per eventuali aggiustamenti futuri.'\}} \\
            }
        }
    }
     \\ \\

    \theutterance \stepcounter{utterance}  
    & & \multicolumn{4}{p{0.6\linewidth}}{
        \cellcolor[rgb]{0.9,0.9,0.9}{
            \makecell[{{p{\linewidth}}}]{
                \texttt{\tiny{[P1$\langle$GM]}}
                \texttt{RIFIUTO: \{'B33', 'C08', 'B20', 'B31', 'A11', 'A82', 'B96', 'B09', 'B93', 'B21'\}} \\
\\ 
\texttt{PROPOSTA: \{'B20', 'B96', 'C08', 'A82', 'B31', 'A21', 'C85', 'B39'\}} \\
\\ 
\texttt{ARGOMENTAZIONE: \{'La tua proposta include alcuni oggetti con efficienza molto bassa rispetto allo sforzo richiesto. La mia proposta ottimizza meglio il rapporto valore{-}sforzo, utilizzando 2991 unità di sforzo su 3357 disponibili. Ho selezionato oggetti strategicamente importanti mantenendo un buon margine di sicurezza di 366 unità per eventuali aggiustamenti futuri.'\}} \\
            }
        }
    }
    & \\ \\

    \theutterance \stepcounter{utterance}  
    & \multicolumn{4}{p{0.6\linewidth}}{
        \cellcolor[rgb]{0.9,0.9,0.9}{
            \makecell[{{p{\linewidth}}}]{
                \texttt{\tiny{[P1$\rangle$GM]}}
                \texttt{RIFIUTO: \{'B20', 'B96', 'C08', 'A82', 'B31', 'A21', 'C85', 'B39'\}} \\
\\ 
\texttt{PROPOSTA: \{'B33', 'C08', 'B20', 'B31', 'A11', 'A82', 'B96', 'B09'\}} \\
\\ 
\texttt{ARGOMENTAZIONE: \{'La tua proposta include oggetti con scarsa efficienza come B39 e C85 che richiedono sforzo elevato. Ho rimosso gli oggetti meno efficienti dalla mia proposta precedente per creare una soluzione più bilanciata. Questa nuova proposta utilizza 2364 unità di sforzo, ben sotto il limite, offrendo una combinazione ottimale che dovrebbe soddisfare entrambi i nostri interessi strategici.'\}} \\
            }
        }
    }
    & & \\ \\

    \theutterance \stepcounter{utterance}  
    & & \multicolumn{4}{p{0.6\linewidth}}{
        \cellcolor[rgb]{0.9,0.9,0.9}{
            \makecell[{{p{\linewidth}}}]{
                \texttt{\tiny{[GM$\rangle$P2]}}
                \texttt{RIFIUTO: \{'B20', 'B96', 'C08', 'A82', 'B31', 'A21', 'C85', 'B39'\}} \\
\\ 
\texttt{PROPOSTA: \{'B33', 'C08', 'B20', 'B31', 'A11', 'A82', 'B96', 'B09'\}} \\
\\ 
\texttt{ARGOMENTAZIONE: \{'La tua proposta include oggetti con scarsa efficienza come B39 e C85 che richiedono sforzo elevato. Ho rimosso gli oggetti meno efficienti dalla mia proposta precedente per creare una soluzione più bilanciata. Questa nuova proposta utilizza 2364 unità di sforzo, ben sotto il limite, offrendo una combinazione ottimale che dovrebbe soddisfare entrambi i nostri interessi strategici.'\}} \\
            }
        }
    }
    & \\ \\

    \theutterance \stepcounter{utterance}  
    & & & \multicolumn{4}{p{0.6\linewidth}}{
        \cellcolor[rgb]{0.9,0.9,0.9}{
            \makecell[{{p{\linewidth}}}]{
                \texttt{\tiny{[GM$\langle$P2]}}
                \texttt{RIFIUTO: \{'B33', 'C08', 'B20', 'B31', 'A11', 'A82', 'B96', 'B09'\}} \\
\\ 
\texttt{PROPOSTA: \{'B33', 'C08', 'B20', 'B31', 'A82', 'B96', 'B09', 'A21', 'C85'\}} \\
\\ 
\texttt{ARGOMENTAZIONE: \{'La tua proposta è sulla strada giusta ma sottoutilizza il nostro budget di sforzo disponibile. Ho mantenuto la maggior parte dei tuoi oggetti ma sostituito A11 con oggetti più strategici come A21 e C85. Questa configurazione utilizza 3100 unità di sforzo rimanendo comunque sotto il limite di 3357 con un margine di sicurezza di 257 unità, massimizzando il valore complessivo del nostro accordo.'\}} \\
            }
        }
    }
     \\ \\

    \theutterance \stepcounter{utterance}  
    & & \multicolumn{4}{p{0.6\linewidth}}{
        \cellcolor[rgb]{0.9,0.9,0.9}{
            \makecell[{{p{\linewidth}}}]{
                \texttt{\tiny{[P1$\langle$GM]}}
                \texttt{RIFIUTO: \{'B33', 'C08', 'B20', 'B31', 'A11', 'A82', 'B96', 'B09'\}} \\
\\ 
\texttt{PROPOSTA: \{'B33', 'C08', 'B20', 'B31', 'A82', 'B96', 'B09', 'A21', 'C85'\}} \\
\\ 
\texttt{ARGOMENTAZIONE: \{'La tua proposta è sulla strada giusta ma sottoutilizza il nostro budget di sforzo disponibile. Ho mantenuto la maggior parte dei tuoi oggetti ma sostituito A11 con oggetti più strategici come A21 e C85. Questa configurazione utilizza 3100 unità di sforzo rimanendo comunque sotto il limite di 3357 con un margine di sicurezza di 257 unità, massimizzando il valore complessivo del nostro accordo.'\}} \\
            }
        }
    }
    & \\ \\

    \theutterance \stepcounter{utterance}  
    & \multicolumn{4}{p{0.6\linewidth}}{
        \cellcolor[rgb]{0.9,0.9,0.9}{
            \makecell[{{p{\linewidth}}}]{
                \texttt{\tiny{[P1$\rangle$GM]}}
                \texttt{RIFIUTO: \{'B33', 'C08', 'B20', 'B31', 'A82', 'B96', 'B09', 'A21', 'C85'\}} \\
\\ 
\texttt{PROPOSTA: \{'B33', 'C08', 'B20', 'B31', 'A11', 'A82', 'B96', 'B09', 'B93'\}} \\
\\ 
\texttt{ARGOMENTAZIONE: \{'Il tuo approccio di sostituire A11 con A21 e C85 non è strategicamente ottimale. A11 offre un rapporto valore{-}sforzo molto superiore rispetto a quegli oggetti. Ho mantenuto la base solida della nostra negoziazione ma aggiunto B93 al posto di A21 e C85, creando una soluzione più efficiente. Questa proposta utilizza 2916 unità di sforzo, lasciando ancora 441 unità di margine, e offre una combinazione più bilanciata per entrambi.'\}} \\
            }
        }
    }
    & & \\ \\

    \theutterance \stepcounter{utterance}  
    & & \multicolumn{4}{p{0.6\linewidth}}{
        \cellcolor[rgb]{0.9,0.9,0.9}{
            \makecell[{{p{\linewidth}}}]{
                \texttt{\tiny{[GM$\rangle$P2]}}
                \texttt{RIFIUTO: \{'B33', 'C08', 'B20', 'B31', 'A82', 'B96', 'B09', 'A21', 'C85'\}} \\
\\ 
\texttt{PROPOSTA: \{'B33', 'C08', 'B20', 'B31', 'A11', 'A82', 'B96', 'B09', 'B93'\}} \\
\\ 
\texttt{ARGOMENTAZIONE: \{'Il tuo approccio di sostituire A11 con A21 e C85 non è strategicamente ottimale. A11 offre un rapporto valore{-}sforzo molto superiore rispetto a quegli oggetti. Ho mantenuto la base solida della nostra negoziazione ma aggiunto B93 al posto di A21 e C85, creando una soluzione più efficiente. Questa proposta utilizza 2916 unità di sforzo, lasciando ancora 441 unità di margine, e offre una combinazione più bilanciata per entrambi.'\}} \\
            }
        }
    }
    & \\ \\

    \theutterance \stepcounter{utterance}  
    & & & \multicolumn{4}{p{0.6\linewidth}}{
        \cellcolor[rgb]{0.9,0.9,0.9}{
            \makecell[{{p{\linewidth}}}]{
                \texttt{\tiny{[GM$\langle$P2]}}
                \texttt{RIFIUTO: \{'B33', 'C08', 'B20', 'B31', 'A11', 'A82', 'B96', 'B09', 'B93'\}} \\
\\ 
\texttt{PROPOSTA: \{'B33', 'C08', 'B20', 'B31', 'A82', 'B96', 'B09', 'B93', 'A21', 'C85'\}} \\
\\ 
\texttt{ARGOMENTAZIONE: \{'Apprezzo la tua proposta ma possiamo ottimizzare ulteriormente utilizzando meglio il nostro budget. Ho sostituito A11 che ha efficienza molto bassa con A21 e aggiunto C85 per massimizzare il valore complessivo. Questa combinazione utilizza 3248 unità di sforzo rimanendo sotto il limite con margine di 109 unità. È una soluzione win{-}win che sfrutta quasi completamente le nostre risorse disponibili.'\}} \\
            }
        }
    }
     \\ \\

    \theutterance \stepcounter{utterance}  
    & & \multicolumn{4}{p{0.6\linewidth}}{
        \cellcolor[rgb]{0.9,0.9,0.9}{
            \makecell[{{p{\linewidth}}}]{
                \texttt{\tiny{[P1$\langle$GM]}}
                \texttt{RIFIUTO: \{'B33', 'C08', 'B20', 'B31', 'A11', 'A82', 'B96', 'B09', 'B93'\}} \\
\\ 
\texttt{PROPOSTA: \{'B33', 'C08', 'B20', 'B31', 'A82', 'B96', 'B09', 'B93', 'A21', 'C85'\}} \\
\\ 
\texttt{ARGOMENTAZIONE: \{'Apprezzo la tua proposta ma possiamo ottimizzare ulteriormente utilizzando meglio il nostro budget. Ho sostituito A11 che ha efficienza molto bassa con A21 e aggiunto C85 per massimizzare il valore complessivo. Questa combinazione utilizza 3248 unità di sforzo rimanendo sotto il limite con margine di 109 unità. È una soluzione win{-}win che sfrutta quasi completamente le nostre risorse disponibili.'\}} \\
            }
        }
    }
    & \\ \\

    \theutterance \stepcounter{utterance}  
    & \multicolumn{4}{p{0.6\linewidth}}{
        \cellcolor[rgb]{0.9,0.9,0.9}{
            \makecell[{{p{\linewidth}}}]{
                \texttt{\tiny{[P1$\rangle$GM]}}
                \texttt{RIFIUTO: \{'B33', 'C08', 'B20', 'B31', 'A82', 'B96', 'B09', 'B93', 'A21', 'C85'\}} \\
\\ 
\texttt{PROPOSTA: \{'B33', 'C08', 'B20', 'B31', 'A11', 'A82', 'B96', 'B09', 'B93', 'B21'\}} \\
\\ 
\texttt{ARGOMENTAZIONE: \{'La tua proposta sacrifica elementi strategicamente cruciali per aggiungere oggetti meno efficienti. A11 rappresenta un asset fondamentale che non dovremmo abbandonare, mentre A21 e C85 offrono un rapporto valore{-}sforzo discutibile. Ho sostituito quei due oggetti con A11 e B21, mantenendo un utilizzo ottimale delle risorse con 3084 unità di sforzo. Questa combinazione offre maggiore stabilità strategica e valore a lungo termine per entrambi.'\}} \\
            }
        }
    }
    & & \\ \\

    \theutterance \stepcounter{utterance}  
    & & \multicolumn{4}{p{0.6\linewidth}}{
        \cellcolor[rgb]{0.9,0.9,0.9}{
            \makecell[{{p{\linewidth}}}]{
                \texttt{\tiny{[GM$\rangle$P2]}}
                \texttt{RIFIUTO: \{'B33', 'C08', 'B20', 'B31', 'A82', 'B96', 'B09', 'B93', 'A21', 'C85'\}} \\
\\ 
\texttt{PROPOSTA: \{'B33', 'C08', 'B20', 'B31', 'A11', 'A82', 'B96', 'B09', 'B93', 'B21'\}} \\
\\ 
\texttt{ARGOMENTAZIONE: \{'La tua proposta sacrifica elementi strategicamente cruciali per aggiungere oggetti meno efficienti. A11 rappresenta un asset fondamentale che non dovremmo abbandonare, mentre A21 e C85 offrono un rapporto valore{-}sforzo discutibile. Ho sostituito quei due oggetti con A11 e B21, mantenendo un utilizzo ottimale delle risorse con 3084 unità di sforzo. Questa combinazione offre maggiore stabilità strategica e valore a lungo termine per entrambi.'\}} \\
            }
        }
    }
    & \\ \\

    \theutterance \stepcounter{utterance}  
    & & & \multicolumn{4}{p{0.6\linewidth}}{
        \cellcolor[rgb]{0.9,0.9,0.9}{
            \makecell[{{p{\linewidth}}}]{
                \texttt{\tiny{[GM$\langle$P2]}}
                \texttt{RIFIUTO: \{'B33', 'C08', 'B20', 'B31', 'A11', 'A82', 'B96', 'B09', 'B93', 'B21'\}} \\
\\ 
\texttt{PROPOSTA: \{'B33', 'C08', 'B20', 'B31', 'A82', 'B96', 'B09', 'B93', 'A21', 'B21'\}} \\
\\ 
\texttt{ARGOMENTAZIONE: \{'La tua insistenza su A11 non è giustificata dai numeri. Ho trovato una soluzione di compromesso che mantiene la maggior parte degli oggetti che consideri importanti ma sostituisce A11 con A21 che offre prestazioni superiori. Questa proposta utilizza 3162 unità di sforzo, rimanendo sotto il limite con margine di 195 unità. È una soluzione equilibrata che dovrebbe soddisfare entrambe le nostre priorità strategiche mantenendo alta efficienza complessiva.'\}} \\
            }
        }
    }
     \\ \\

    \theutterance \stepcounter{utterance}  
    & & \multicolumn{4}{p{0.6\linewidth}}{
        \cellcolor[rgb]{0.9,0.9,0.9}{
            \makecell[{{p{\linewidth}}}]{
                \texttt{\tiny{[P1$\langle$GM]}}
                \texttt{RIFIUTO: \{'B33', 'C08', 'B20', 'B31', 'A11', 'A82', 'B96', 'B09', 'B93', 'B21'\}} \\
\\ 
\texttt{PROPOSTA: \{'B33', 'C08', 'B20', 'B31', 'A82', 'B96', 'B09', 'B93', 'A21', 'B21'\}} \\
\\ 
\texttt{ARGOMENTAZIONE: \{'La tua insistenza su A11 non è giustificata dai numeri. Ho trovato una soluzione di compromesso che mantiene la maggior parte degli oggetti che consideri importanti ma sostituisce A11 con A21 che offre prestazioni superiori. Questa proposta utilizza 3162 unità di sforzo, rimanendo sotto il limite con margine di 195 unità. È una soluzione equilibrata che dovrebbe soddisfare entrambe le nostre priorità strategiche mantenendo alta efficienza complessiva.'\}} \\
            }
        }
    }
    & \\ \\

    \theutterance \stepcounter{utterance}  
    & \multicolumn{4}{p{0.6\linewidth}}{
        \cellcolor[rgb]{0.9,0.9,0.9}{
            \makecell[{{p{\linewidth}}}]{
                \texttt{\tiny{[P1$\rangle$GM]}}
                \texttt{RIFIUTO: \{'B33', 'C08', 'B20', 'B31', 'A82', 'B96', 'B09', 'B93', 'A21', 'B21'\}} \\
\\ 
\texttt{PROPOSTA: \{'B33', 'C08', 'B20', 'B31', 'A11', 'B96', 'B09', 'A21', 'B21'\}} \\
\\ 
\texttt{ARGOMENTAZIONE: \{'Ho trovato una soluzione di compromesso che include sia A11 che A21 che sembrano essere priorità per entrambi. Ho rimosso A82 e B93 per fare spazio a entrambi gli oggetti, utilizzando 3318 unità di sforzo rimanendo sotto il limite con margine di 39 unità. Questa proposta rispetta le tue preferenze per A21 mantenendo A11 che considero strategicamente essenziale, creando una vera soluzione win{-}win.'\}} \\
            }
        }
    }
    & & \\ \\

    \theutterance \stepcounter{utterance}  
    & & \multicolumn{4}{p{0.6\linewidth}}{
        \cellcolor[rgb]{0.9,0.9,0.9}{
            \makecell[{{p{\linewidth}}}]{
                \texttt{\tiny{[GM$\rangle$P2]}}
                \texttt{RIFIUTO: \{'B33', 'C08', 'B20', 'B31', 'A82', 'B96', 'B09', 'B93', 'A21', 'B21'\}} \\
\\ 
\texttt{PROPOSTA: \{'B33', 'C08', 'B20', 'B31', 'A11', 'B96', 'B09', 'A21', 'B21'\}} \\
\\ 
\texttt{ARGOMENTAZIONE: \{'Ho trovato una soluzione di compromesso che include sia A11 che A21 che sembrano essere priorità per entrambi. Ho rimosso A82 e B93 per fare spazio a entrambi gli oggetti, utilizzando 3318 unità di sforzo rimanendo sotto il limite con margine di 39 unità. Questa proposta rispetta le tue preferenze per A21 mantenendo A11 che considero strategicamente essenziale, creando una vera soluzione win{-}win.'\}} \\
            }
        }
    }
    & \\ \\

    \theutterance \stepcounter{utterance}  
    & & & \multicolumn{4}{p{0.6\linewidth}}{
        \cellcolor[rgb]{0.9,0.9,0.9}{
            \makecell[{{p{\linewidth}}}]{
                \texttt{\tiny{[GM$\langle$P2]}}
                \texttt{RIFIUTO: \{'B33', 'C08', 'B20', 'B31', 'A11', 'B96', 'B09', 'A21', 'B21'\}} \\
\\ 
\texttt{PROPOSTA: \{'B33', 'C08', 'B20', 'B31', 'A82', 'B96', 'B09', 'A21', 'B21'\}} \\
\\ 
\texttt{ARGOMENTAZIONE: \{'La tua proposta è interessante ma A82 offre caratteristiche tecniche superiori rispetto ad A11 per il nostro progetto complessivo. Ho sostituito A11 con A82 mantenendo tutti gli altri oggetti che hai proposto, inclusi sia A21 che B21. Questa configurazione utilizza 3162 unità di sforzo con margine di 195 unità, offrendo prestazioni ottimali e stabilità per entrambi. È una soluzione che combina i migliori elementi delle nostre proposte precedenti.'\}} \\
            }
        }
    }
     \\ \\

    \theutterance \stepcounter{utterance}  
    & & \multicolumn{4}{p{0.6\linewidth}}{
        \cellcolor[rgb]{0.9,0.9,0.9}{
            \makecell[{{p{\linewidth}}}]{
                \texttt{\tiny{[P1$\langle$GM]}}
                \texttt{RIFIUTO: \{'B33', 'C08', 'B20', 'B31', 'A11', 'B96', 'B09', 'A21', 'B21'\}} \\
\\ 
\texttt{PROPOSTA: \{'B33', 'C08', 'B20', 'B31', 'A82', 'B96', 'B09', 'A21', 'B21'\}} \\
\\ 
\texttt{ARGOMENTAZIONE: \{'La tua proposta è interessante ma A82 offre caratteristiche tecniche superiori rispetto ad A11 per il nostro progetto complessivo. Ho sostituito A11 con A82 mantenendo tutti gli altri oggetti che hai proposto, inclusi sia A21 che B21. Questa configurazione utilizza 3162 unità di sforzo con margine di 195 unità, offrendo prestazioni ottimali e stabilità per entrambi. È una soluzione che combina i migliori elementi delle nostre proposte precedenti.'\}} \\
            }
        }
    }
    & \\ \\

    \theutterance \stepcounter{utterance}  
    & \multicolumn{4}{p{0.6\linewidth}}{
        \cellcolor[rgb]{0.9,0.9,0.9}{
            \makecell[{{p{\linewidth}}}]{
                \texttt{\tiny{[P1$\rangle$GM]}}
                \texttt{RIFIUTO: \{'B33', 'C08', 'B20', 'B31', 'A82', 'B96', 'B09', 'A21', 'B21'\}} \\
\\ 
\texttt{PROPOSTA: \{'B33', 'C08', 'B20', 'B31', 'A11', 'A82', 'B96', 'B09', 'B21'\}} \\
\\ 
\texttt{ARGOMENTAZIONE: \{'La tua proposta ha merit ma possiamo ottimizzare ulteriormente includendo sia A11 che A82 per massimizzare le capacità tecniche complessive. Ho rimosso solo A21 per fare spazio ad entrambi gli oggetti tecnici principali. Questa configurazione utilizza 2812 unità di sforzo, lasciando un margine confortevole di 545 unità, e combina i punti di forza di entrambe le nostre visioni strategiche per una soluzione più robusta e versatile.'\}} \\
            }
        }
    }
    & & \\ \\

    \theutterance \stepcounter{utterance}  
    & & \multicolumn{4}{p{0.6\linewidth}}{
        \cellcolor[rgb]{0.9,0.9,0.9}{
            \makecell[{{p{\linewidth}}}]{
                \texttt{\tiny{[GM$\rangle$P2]}}
                \texttt{RIFIUTO: \{'B33', 'C08', 'B20', 'B31', 'A82', 'B96', 'B09', 'A21', 'B21'\}} \\
\\ 
\texttt{PROPOSTA: \{'B33', 'C08', 'B20', 'B31', 'A11', 'A82', 'B96', 'B09', 'B21'\}} \\
\\ 
\texttt{ARGOMENTAZIONE: \{'La tua proposta ha merit ma possiamo ottimizzare ulteriormente includendo sia A11 che A82 per massimizzare le capacità tecniche complessive. Ho rimosso solo A21 per fare spazio ad entrambi gli oggetti tecnici principali. Questa configurazione utilizza 2812 unità di sforzo, lasciando un margine confortevole di 545 unità, e combina i punti di forza di entrambe le nostre visioni strategiche per una soluzione più robusta e versatile.'\}} \\
            }
        }
    }
    & \\ \\

    \theutterance \stepcounter{utterance}  
    & & & \multicolumn{4}{p{0.6\linewidth}}{
        \cellcolor[rgb]{0.9,0.9,0.9}{
            \makecell[{{p{\linewidth}}}]{
                \texttt{\tiny{[GM$\langle$P2]}}
                \texttt{RIFIUTO: \{'B33', 'C08', 'B20', 'B31', 'A11', 'A82', 'B96', 'B09', 'B21'\}} \\
\\ 
\texttt{PROPOSTA: \{'B33', 'C08', 'B20', 'B31', 'A82', 'B96', 'B09', 'A21', 'B21'\}} \\
\\ 
\texttt{ARGOMENTAZIONE: \{'La tua proposta è ragionevole ma sottoutilizza il nostro budget disponibile. Abbiamo 3357 unità a disposizione e la tua proposta ne usa solo 2592, sprecando 765 unità di capacità. A21 è un asset strategico fondamentale che non possiamo permetterci di escludere {-} offre caratteristiche uniche essenziali per il successo del progetto. Ho sostituito A11 con A21, utilizzando 3162 unità e lasciando ancora margine adeguato di 195 unità.'\}} \\
            }
        }
    }
     \\ \\

    \theutterance \stepcounter{utterance}  
    & & \multicolumn{4}{p{0.6\linewidth}}{
        \cellcolor[rgb]{0.9,0.9,0.9}{
            \makecell[{{p{\linewidth}}}]{
                \texttt{\tiny{[P1$\langle$GM]}}
                \texttt{RIFIUTO: \{'B33', 'C08', 'B20', 'B31', 'A11', 'A82', 'B96', 'B09', 'B21'\}} \\
\\ 
\texttt{PROPOSTA: \{'B33', 'C08', 'B20', 'B31', 'A82', 'B96', 'B09', 'A21', 'B21'\}} \\
\\ 
\texttt{ARGOMENTAZIONE: \{'La tua proposta è ragionevole ma sottoutilizza il nostro budget disponibile. Abbiamo 3357 unità a disposizione e la tua proposta ne usa solo 2592, sprecando 765 unità di capacità. A21 è un asset strategico fondamentale che non possiamo permetterci di escludere {-} offre caratteristiche uniche essenziali per il successo del progetto. Ho sostituito A11 con A21, utilizzando 3162 unità e lasciando ancora margine adeguato di 195 unità.'\}} \\
            }
        }
    }
    & \\ \\

    \theutterance \stepcounter{utterance}  
    & \multicolumn{4}{p{0.6\linewidth}}{
        \cellcolor[rgb]{0.9,0.9,0.9}{
            \makecell[{{p{\linewidth}}}]{
                \texttt{\tiny{[P1$\rangle$GM]}}
                \texttt{ACCORDO: \{'B33', 'C08', 'B20', 'B31', 'A82', 'B96', 'B09', 'A21', 'B21'\}} \\
\\ 
\texttt{ARGOMENTAZIONE: \{'Dopo attenta considerazione, riconosco che la tua proposta rappresenta un buon equilibrio tra le nostre priorità diverse. Anche se personalmente preferivo mantenere A11, capisco che A21 ha caratteristiche specifiche che ritieni fondamentali per il successo del progetto. La configurazione proposta utilizza efficacemente le nostre risorse disponibili e dovrebbe fornire risultati soddisfacenti per entrambi. Accetto questa soluzione come base per procedere.'\}} \\
            }
        }
    }
    & & \\ \\

    \theutterance \stepcounter{utterance}  
    & & & \multicolumn{2}{p{0.3\linewidth}}{
        \cellcolor[rgb]{0.9,0.9,0.9}{
            \makecell[{{p{\linewidth}}}]{
                \texttt{\tiny{[GM$|$GM]}}
                \texttt{game successful} \\
            }
        }
    }
    & & \\ \\

    \theutterance \stepcounter{utterance}  
    & & & \multicolumn{2}{p{0.3\linewidth}}{
        \cellcolor[rgb]{0.9,0.9,0.9}{
            \makecell[{{p{\linewidth}}}]{
                \texttt{\tiny{[GM$|$GM]}}
                \texttt{end game} \\
            }
        }
    }
    & & \\ \\

\end{supertabular}
}

\end{document}
